This section reports aggregate performance, cross-competition consistency, reliability and stability, scaling with time budget, and per-competition highlights. The headline performance leaderboard uses a normalization that enables comparisons across competitions and time budgets.

\subsection{Key findings}
\begin{itemize}
  \item Under the headline normalization (best budget per competition, averaged across competitions), \texttt{MiniMax-M2.1-TEE} is the top performer and is rank-1 in every competition.
  \item Reliability varies meaningfully even among strong performers. Success-rate and stability plots show clear separation between more and less reliable models.
  \item Some models gain substantially from extra time budget, while others are relatively flat. Marginal-gain and monotonicity views capture these patterns.
\end{itemize}

\subsection{Aggregate performance leaderboard (headline)}
The aggregate leaderboard is derived from five-run medians and normalized via rank-points so that scores from different competitions (AUC vs RMSE) are comparable.

Method:
\begin{itemize}
  \item Unit of aggregation: per \texttt{(competition, budget)} setting, use the model's five-run median \texttt{score\_raw}.
  \item Normalize within each setting by rank: best model gets 1.0, worst gets 0.0, with linear spacing in between.
  \item Headline aggregation: for each \texttt{(model, competition)} take the best normalized setting across the three budgets, then average across the 4 competitions with equal weights.
\end{itemize}

This rank-based normalization is applied after accounting for metric direction (for example, AUC is higher-is-better, while RMSE is lower-is-better).

Interpretation: rank-points are relative within each \texttt{(competition, budget)} setting. They preserve ordering within a setting rather than absolute metric gaps.

\begin{figure}[t]
  \centering
  \includegraphics[width=\linewidth]{leaderboard_headline_best_budget.png}
  \caption{Headline leaderboard: best budget per competition.}
\end{figure}

Robustness variants (secondary) include: (i) an overall aggregation that averages all \texttt{(competition, budget)} settings equally, and (ii) a 1200s-only aggregation. See Appendix~\ref{sec:appendix-robustness}.

\subsection{Cross-competition consistency}
Per-competition ranks are computed in the same normalized space as the headline leaderboard (best budget per competition). The heatmap below shows each model's rank (1=best) per competition.

\begin{figure}[t]
  \centering
  \includegraphics[width=\linewidth]{consistency_ranks_heatmap.png}
  \caption{Per-competition ranks (1=best).}
\end{figure}

Rank variability across competitions is summarized via rank standard deviation (lower is more consistent). See Appendix~\ref{sec:appendix-consistency}.

\subsection{Reliability and stability}
Reliability has two components: run success rate (how often a run yields a valid score) and within-setting stability (how variable a model is across the five runs used for each reported setting). The trade-off is summarized via a Pareto-style plot (performance vs stability; color indicates success rate).

\begin{figure}[t]
  \centering
  \includegraphics[width=\linewidth]{reliability_pareto_performance_vs_stability.png}
  \caption{Performance vs stability (color indicates success rate).}
\end{figure}

Supporting breakdown plots for success rate and stability are included in Appendix~\ref{sec:appendix-reliability}.

\subsection{Scaling with time budget}
Normalized performance is analyzed as time budget increases from 240s to 600s to 1200s, averaged across the four competitions.

On aggregate, scaling is broadly consistent with the expected monotonic pattern. Define a model's per-competition scaling curve as monotone if its direction-corrected five-run medians do not worsen as budget increases (240s $\rightarrow$ 600s $\rightarrow$ 1200s). Under this definition: (i) across all 40 model$\times$competition curves, 23/40 (57.5\%) are monotone, and (ii) across the 10 models, the median model is monotone in 62.5\% of competitions. Appendix~\ref{sec:appendix-scaling} reports monotonicity rates and marginal gains.

\begin{figure}[t]
  \centering
  \includegraphics[width=\linewidth]{scaling_points_lines.png}
  \caption{Scaling with time budget.}
\end{figure}

At the individual model level, scaling can be noisy. Each task$\times$budget setting is summarized by five successful runs. More runs are likely required for stable model-level scaling curves. Appendix~\ref{sec:appendix-scaling} reports marginal gains and monotonicity rates.

\subsection{Per-competition highlights}
This subsection highlights a small number of representative task/budget settings. Full five-run tables are available in the companion repository materials.

\paragraph{bank-customer-churn-ict-u-ai (AUC; higher is better).}
The strongest median AUC at 1200s is 0.928000 (GPT OSS 120B TEE).
At 240s, the top median is 0.926671 (MiniMax-M2.1-TEE).
This task shows top-tier clustering in the 0.92x range, with notable underperformance from some models in specific settings (for example, NVIDIA-Nemotron-3-Nano at 0.813105 at 1200s).

\paragraph{foot-traffic-wuerzburg-retail-forecasting-2-0 (RMSE; lower is better).}
MiniMax-M2.1-TEE is best at all three budgets (0.066846, 0.065770, 0.065489).
A key instability signal appears for GLM 4.7 Flash at 1200s: median 0.107502 with IQR 0.070186..0.221725. This IQR is substantially wider than neighboring models.

\paragraph{playground-series-s5e10 (RMSE; lower is better).}
At 1200s, medians are tightly clustered near 0.0562, with the best setting at 0.056190 (GLM-4.6-FP8) and many models within a few $1\mathrm{e}{-4}$.

\paragraph{playground-series-s6e1 (RMSE; lower is better).}
MiniMax-M2.1-TEE leads at 1200s with RMSE 8.699779.
At 600s, TNG-R1T2-Chimera has a large failure-mode outlier (median 10.199380, IQR 9.088197..13.444163). This motivates caution when interpreting single-budget standings.
